\homework{2}{Fri 10 Oct 2025 23:16}{}

\begin{itemize}
	\item
		2.1. a) Show that the conclusion of Sanov's Theorem doesn't hold as stated in class for all closed sets. That is, show that there exists a closed set $F \subset \mathcal{P}_m$ such that

\[
\lim _{n \rightarrow \infty} \frac{1}{n} \log P\left(L_n \in F\right) \neq-\inf _{\nu \in F} H(\nu \mid \rho) .
\]

b) Show that for any closed subset $F \subset \mathcal{P}_m$ we have

\[
\limsup _{n \rightarrow \infty} \frac{1}{n} \log P\left(L_n \in F\right) \leq-\inf _{\nu \in F} H(\nu \mid \rho) .
\]


Hint: Consider open subsets $G^{\varepsilon} \searrow F$ as $\varepsilon \rightarrow 0$.
\begin{proof}
	Pick $F = \{\nu\} $ where $\nu \in \mathcal{P}_m$ corresponds to an irrational point in the probability simplex. In other words, $\nu = \sum_j \nu_j \delta_j$ where $\nu_j \in [0,1]\setminus \mathbb{Q}$ for at least one $j$. Then, we have $L_n \not\in F$ for all $n$.  Thus the left hand side $\frac{1}{n} \log P(L_n \in F) = - \infty $ and the limit does not exist.

	Now, for any closed subset $F \subset \mathcal{P}_m$, we can pick $G^\varepsilon$ be $\varepsilon$-exterior of $F$, so $F \subset G^\varepsilon$ and $G^\varepsilon$ open.
	Then
	\[
		P(L_n \in F) \le P(L_n \in G^\varepsilon)
	.\] 
	So that
	\[
		\limsup \frac{1}{n} \log P(L_n \in F) \le 
		\limsup \frac{1}{n} \log P(L_n \in G^\varepsilon) =
		- \inf_{\nu \in G^\varepsilon} H(\nu | \rho)
	.\] 
	This is true for any $\varepsilon > 0$, so
	\[
		\limsup \frac{1}{n} \log P(L_n \in F) \le 
		- \sup  _{\varepsilon > 0} \inf _{\nu \in G^\varepsilon} H(\nu | \rho)
	.\] 
	Since the family $G^\varepsilon$ is monotone decreasing as $\varepsilon \to 0$ and $\cap_{\varepsilon > 0} G^\varepsilon = F$,
	\[
		\sup  _{\varepsilon > 0} \inf _{\nu \in G^\varepsilon} H(\nu | \rho) = \inf _{\nu \in F} H(\nu | \rho)
	.\] 
\end{proof}

\item 
	Show the identity claimed in class for the large deviation rate function $I_\rho^2(\nu)$ of the pair empirical measure $L_n^2=\frac{1}{n} \sum_{i=1}^n \delta_{\left(X_i, X_{i+1}\right)}$. That is, if

\[
I_\rho^2(\nu)=\sum_{i, j} \nu_{i, j} \log \left(\frac{\nu_{i, j}}{\bar{\nu}_i \rho_j}\right)
\]

then show that

\[
I_\rho^2(\nu)=H(\bar{\nu} \mid \rho)+H(\nu \mid \bar{\nu} \times \bar{\nu})
\]
\begin{proof}
	\begin{align*}
		& H(\overline{\nu}  | \rho) + H(\nu | \overline{\nu}  \times \overline{\nu} )\\
		&= \sum_i \overline{\nu} _i \log \frac{\overline{\nu} _i}{\rho_i} + \sum_{i j}\nu_{i j}\log \frac{\nu_{i j}}{\overline{\nu} _i \overline{\nu} _j} \\
		&= \sum_{i j} \nu_{i j}\left( \log \frac{\overline{\nu} _i}{\rho_i} + \log \frac{\nu_{i j}}{\overline{\nu} _i \overline{\nu} _j} \right)  \\
		&= \sum_{i j} \nu_{i j} \log \frac{\nu_{i j}}{\rho_i \overline{\nu} _j} \\
		&= I^2_\rho(\nu)
	.\end{align*}
\end{proof}

\item
The relative entropy function $H(\nu \mid \rho)$ is defined by

\[
H(\nu \mid \rho)= \begin{cases}\int \frac{d \nu}{d \rho} \log \left(\frac{d \nu}{d \rho}\right) d \rho & \text { if } \nu \ll \rho \\ \infty & \text { otherwise }\end{cases}
\]


Notes:


- Recall that $\nu \ll \rho$ means that $\nu$ is absolutely continuous with respect to $\rho$ and so the Radon-Nikodym derivative $\frac{d \nu}{d \rho}$ exists.

- In class we've mostly been dealing with the relative entropy in the setting of probability measures on finite state spaces, but for this problem we are in the more general setting.

a) Prove that the relative entropy function $\nu \mapsto H(\nu \mid \rho)$ is strictly convex (partial credit if you only prove convexity) on the domain where it is finite. That is, if $\nu \ll \rho, \mu \ll \rho$ and $\nu \neq \mu$ then
\[
H(t \nu+(1-t) \mu \mid \rho)<t H(\nu \mid \rho)+(1-t) H(\mu \mid \rho), \quad \forall t \in(0,1) .
\]
\begin{proof}
	Write $f = \frac{d \mu}{d \rho}$ and $g = \frac{d \nu}{d \rho}$. Suppressing $d \rho$ from the integral notation, we would like to prove that
	\[
		\int t f \log f + (1-t) g \log g - (t f + (1 - t) g) \log \left( t f + (1 - t) g \right) >  0
	.\] 
	Now, since $x \to x \log  x$ is strictly convex, we obtain convexity by integrating the inequality
	\[
		t f(\omega) \log f(\omega) + (1 - t) g(\omega) \log  g(\omega) - (t f (\omega) + (1 - t) g(\omega) \log \left( t f(\omega) + (1 - t) g(\omega) \right)  > 0
	.\] 
	Now, to prove strict convexity, suppose $f - g > \varepsilon$ over a measurable set $A$. Pointwise application of the convexity of $x \mapsto x \log  x$ tells us that there exists some $\eta > 0$ such that
	\[
		t f(\omega) \log f(\omega) + (1 - t) g(\omega) \log  g(\omega) - (t f (\omega) + (1 - t) g(\omega) \log \left( t f(\omega) + (1 - t) g(\omega) \right)  > \eta \qquad
		\text{on }A
	.\] 
	Combining with convexity, this tells us that our integral
	\[
		\int \left( \ldots \right) \ge \eta |A| > 0
	.\] 
\end{proof}

b) Prove that $H(\nu \mid \rho) \geq 0$, and conclude that $H(\nu \mid \rho)=0$ if and only if $\nu=\rho$.
\begin{proof}
	Write $f = \frac{d \nu}{ d \rho}$ which is a measurable function. Approximate $f$ with simple functions $f_n$ from below such that $f_n = \sum_{m=1}^{M_n} c_{n m} 1_{A_{n m}}$ where $c_{n m}\ge 0$ and $(A_{n m})_{1 \le m \le M_n}$ partitions the probability space $\Omega$. We have
	\[
		\int |f - f_n| \to 0 \text{  as  } n \to \infty 
	.\] 

	Since $f$ is a probability density, $\int  f = 1$, so that for large enough $n$,
	\[
		\int  f_n = \sum_{m = 1}^{M_n} c_{n m}|A_{n m}| \le 1	.\] 
	Therefore each $c_{n m} |A_{n m}| \le  1$. Now
	\begin{align*}
		&\int f_n \log f_n\\
		&= \sum_m |A_{n m}| c_{n m} \log c_{n m} \\
		&= \sum_m |A_{n m}| c_{n m} \log (|A_{n m}| c_{n m}) - c_{n m} |A_{n m}| \log |A_{n m}| \\
		&\ge \sum_m - c_{n m} |A_{n m}| \log |A_{n m}| \\
		&\ge 0
	.\end{align*}
	In the first inequality, we used the fact that $[0,1]\ni x \mapsto x \log x \le 0$. In the second inequality, we used the fact that $\sum_m |A_{m n}| = 1$ hence $\log |A_{m n}| \le  0$.

	Finally, application of monotone convergence theorem tells us that
	\[
		\int f_n \log  f_n \to \int  f \log  f \ge 0
	.\] 

	If $H(\nu, \rho) = \int f \log  f = 0$, then $\log f \equiv 0$, so $\frac{d \nu}{d \rho} = 1$, hence $\nu = \rho$ are equivalent measures.
\end{proof}
\end{itemize}

